% chktex-file 9
% chktex-file 17

本研究以蒙地卡羅模擬計算理想條件下（不含雜散光與量測雜訊）的角度分布，作為實驗對照。每顆光子自比色槽入口射入，先依
\[
s=-\frac{\ln\xi}{\mu_s}
\]
前進隨機距離，其中 $\xi$ 為 $0$ 至 $1$ 之間的亂數，$\mu_s$ 為散射係數；碰到微球後，依含粒徑分布的 Mie 相位函數決定新方向，粒徑分布的變異係數分別採用 $0.492\,\mu\mathrm{m}$ 的 $2.1\%$ 與 $2.16\,\mu\mathrm{m}$ 的 $3.1\%$。光子碰到比色槽壁時，依 Fresnel 公式處理水、PMMA（折射率 $1.489$）與空氣界面的反射及透射，包含槽壁內多次反射與全反射；穿出後再依 Snell 定律折射，記錄其落入旋轉臂的 $1^\circ$ 角度格。

每個 $\tau$ 模擬 $5{,}000$ 萬顆光子；光子偏離光束平面超過 $150\,\mathrm{mm}$ 後停止追蹤。模擬假設雷射為非偏振光，旋轉臂半徑取實測值 $110\,\mathrm{mm}$，並以遮光管收光半角 $15.3^\circ$（管長 $11\,\mathrm{mm}$、開口寬 $6\,\mathrm{mm}$）限制感測器收光。每個條件使用 $6$ 組不同亂數種子，取模擬平均值，再以和實驗相同的方式計算 $C$。驗證時，移除比色槽、令 $\tau\to0$ 並只計入單次散射光子，模擬結果與 Mie 理論一致；另以獨立重寫的程式計算六個條件，結果在統計誤差內一致。

模擬方式：
\subsection*{(一) 前置準備}
\begin{itemize}
    \item[1.] 計算出粒子的 $S_1(\theta),\ S_2(\theta)$ 存成查表
    \item[2.] 計算出粒子的散射、吸收截面$\left(\sigma_s,\ \sigma_a\right)$
    \item[3.] 計算出粒子的散射、吸收係數$\left(\mu_s,\ \mu_a\right)$
    \item[4.] 將穆勒矩陣的 $S_{11}$ 轉成累積分布函數 $CDF$ 
\end{itemize}

每輪模擬分別追蹤一次光子進入到容器後消失或離開容器的生存週期，每個光子儲存以下資訊：
\begin{minted}[frame=lines, linenos, fontsize=\small]{rust}
#[repr(align(32))] // 加速執行用的記憶體對齊
pub(crate) struct Photon {
    pub(crate) stokes_vector: StokesVector, // 光子目前的斯托克向量
    pub(crate) direction: DirecVec, // 光子目前的方向 (必為單位向量)
    pub(crate) start_location: DirecVec, // 移動前位置向量
    pub(crate) last_plane_normal_v: DirecVec, // 上一個散射平面的法向量
    pub(crate) w: f64, // 粒子權重
    pub(crate) scattered: bool, // 是否經過散射
}
\end{minted}

\subsection*{(二) 模擬內容}
\subsection*{1. 初始化光子}
初始化尚未射入方形比色管的光子

\begin{minted}[frame=lines, linenos, fontsize=\small]{rust}
pub(crate) const PHOTON: photon::Photon = photon::Photon {
    stokes_vector: Vec4::new(1.0, 0.0, 0.0, 0.0), // 光強為 1，偏振視雷射測量而定
    direction: Vec3::new(1.0, 0.0, 0.0), // 朝 x+ 方向射入方形比色管
    start_location: Vec3::new(0.0, 0.0, 0.0), // 從原點射向比色管
    last_plane_normal_v: Vec3::new(0.0, 0.0, 1.0), // 處使值不會用到，預設水平平面
    w: 1.0, // 權重為一
    scattered: false, // 尚未散射
};
\end{minted}

\subsection*{2. 抽取自由路徑}
\begin{itemize}
    \item[(1)] 按照自由路徑分布公式\eqref{eq:mean_free_path}，隨機出一變數$\zeta$，後藉此找出此次的預計移動長度
    \item[(2)] 使用粒子儲存的移動方向找出下個點的預計位置。
\end{itemize}

\subsection*{3. 確認是否仍在邊界中}
\noindent 是：進入碰撞事件\\
否：使用 Fresnel 公式隨機出是否會離開邊界，如果會，則使用斯乃爾定律求出離開後的角度並記錄下來，如果會再次反射進入比色管，則繼續移動直到路徑剩餘為 0

\subsection*{4. 碰撞事件}
當光子遇到溶液中粒子後則會進入此事件
\begin{itemize}
    \item[(1)] 將粒子銓重扣除 $\frac{\mu_a}{\mu_t}$
    \item[(2)] 使用前置作業已儲存的分布函數 $CDF$ 隨機出轉向角 $\theta$
    \item[(3)] 隨機抽出旋轉角 $\psi [0, 2\pi)$
    \item[(4)] 藉隨機出的角度找出新的方向向量
    \item[(5)] 使用式~\eqref{eq:monte_} 找出新的斯托克向量
\end{itemize}

\subsection*{5. 權重剪枝檢查}
根據自己所定義的最低權重(我們使用 $\frac{1}{10}$)，若該光子權重已小於則做出類似俄羅斯輪盤的隨機死亡，若存活則重新步驟 2
